% Options for packages loaded elsewhere
\PassOptionsToPackage{unicode}{hyperref}
\PassOptionsToPackage{hyphens}{url}
\PassOptionsToPackage{dvipsnames,svgnames,x11names}{xcolor}
\documentclass[
  american,
  11pt,
  a4paper,
]{article}
\usepackage{xcolor}
\usepackage[margin=24mm]{geometry}
\usepackage{amsmath,amssymb}
\setcounter{secnumdepth}{5}
\usepackage{iftex}
\ifPDFTeX
  \usepackage[T1]{fontenc}
  \usepackage[utf8]{inputenc}
  \usepackage{textcomp} % provide euro and other symbols
\else % if luatex or xetex
  \usepackage{unicode-math} % this also loads fontspec
  \defaultfontfeatures{Scale=MatchLowercase}
  \defaultfontfeatures[\rmfamily]{Ligatures=TeX,Scale=1}
\fi
\usepackage{lmodern}
\ifPDFTeX\else
  % xetex/luatex font selection
\fi
% Use upquote if available, for straight quotes in verbatim environments
\IfFileExists{upquote.sty}{\usepackage{upquote}}{}
\IfFileExists{microtype.sty}{% use microtype if available
  \usepackage[]{microtype}
  \UseMicrotypeSet[protrusion]{basicmath} % disable protrusion for tt fonts
}{}
\makeatletter
\@ifundefined{KOMAClassName}{% if non-KOMA class
  \IfFileExists{parskip.sty}{%
    \usepackage{parskip}
  }{% else
    \setlength{\parindent}{0pt}
    \setlength{\parskip}{6pt plus 2pt minus 1pt}}
}{% if KOMA class
  \KOMAoptions{parskip=half}}
\makeatother
\usepackage{longtable,booktabs,array}
\usepackage{calc} % for calculating minipage widths
% Correct order of tables after \paragraph or \subparagraph
\usepackage{etoolbox}
\makeatletter
\patchcmd\longtable{\par}{\if@noskipsec\mbox{}\fi\par}{}{}
\makeatother
% Allow footnotes in longtable head/foot
\IfFileExists{footnotehyper.sty}{\usepackage{footnotehyper}}{\usepackage{footnote}}
\makesavenoteenv{longtable}
\ifLuaTeX
\usepackage[bidi=basic,shorthands=off,]{babel}
\else
\usepackage[bidi=default,shorthands=off,]{babel}
\fi
\ifLuaTeX
  \usepackage{selnolig} % disable illegal ligatures
\fi
\setlength{\emergencystretch}{3em} % prevent overfull lines
\providecommand{\tightlist}{%
  \setlength{\itemsep}{0pt}\setlength{\parskip}{0pt}}
\usepackage[T1]{fontenc}
\usepackage{lmodern}
\usepackage{microtype}
\usepackage{xcolor}
\usepackage{fancyhdr}
\usepackage{amsmath,amssymb}
\usepackage{xurl}
\definecolor{qddmink}{HTML}{17324D}
\pagestyle{fancy}
\fancyhf{}
\fancyhead[L]{\small\textcolor{qddmink}{Quad-Dimensional DataMatrix}}
\fancyhead[R]{\small RfD 0.2.0}
\fancyfoot[C]{\thepage}
\setlength{\headheight}{15pt}
\setlength{\emergencystretch}{3em}

\usepackage{bookmark}
\IfFileExists{xurl.sty}{\usepackage{xurl}}{} % add URL line breaks if available
\urlstyle{same}
\hypersetup{
  pdftitle={Quad-Dimensional DataMatrix},
  pdfauthor={João Paulo De Luca Jr; Jorge Luiz Espanhol},
  pdflang={en-US},
  pdfsubject={Request for Discussion: hexadecimal indexed GIF states},
  colorlinks=true,
  linkcolor={qddmink},
  filecolor={Maroon},
  citecolor={Blue},
  urlcolor={qddmink},
  pdfcreator={LaTeX via pandoc}}

\title{Quad-Dimensional DataMatrix}
\usepackage{etoolbox}
\makeatletter
\providecommand{\subtitle}[1]{% add subtitle to \maketitle
  \apptocmd{\@title}{\par {\large #1 \par}}{}{}
}
\makeatother
\subtitle{Request for Discussion: Hexadecimal GIF States and App-Held
Sequences}
\author{João Paulo De Luca Jr \and Jorge Luiz Espanhol}
\date{4 November 2026 · RfD 0.2.0}

\begin{document}
\maketitle
\begin{abstract}
Quad-Dimensional DataMatrix is proposed as a discussion framework for
hexadecimal data represented by indexed GIF states. The working profile
uses a 256 by 256 pixel canvas, 64 by 64 effective modules, up to 64
hexadecimal palette entries, and up to 64 frames. Four logical
coordinates identify horizontal position, vertical position, temporal
state, and dimension. Their maximum Cartesian catalogue contains
16,777,216 candidate tuples. A dense sequence selects one of 64
alternatives at each module and therefore carries up to 192 KiB of gross
payload in a 64-frame cycle, nominally eight seconds at 8 fps. An
app-held sequence supports a second interaction: advance at a chosen
cadence, pause on user selection, and copy the state's hexadecimal code
for a receiving service. The paper develops the counting model,
separates the two timing roles, and examines one-to-one
Unicode-to-hexadecimal mappings as candidates for a future dictionary or
key. A 16-frame case study and a local bilingual viewer accompany the
discussion. Scope and Limitations consolidates the qualifications. The
document invites refinement before a specification is proposed.
\end{abstract}

\section{Purpose and discussion
status}\label{purpose-and-discussion-status}

\textbf{Quad-Dimensional DataMatrix} is a Request for Discussion (RfD)
by \textbf{João Paulo De Luca Jr and Jorge Luiz Espanhol}. It brings
together a four-coordinate catalogue, indexed GIF transport, hexadecimal
values, and app-held state selection.

The immediate objective is to make the proposed objects countable and
testable. A reader should be able to identify what a state contains, how
the app selects it, what the user copies, and which interpretation of a
hexadecimal value is being proposed. A later stage may develop a
specification from the discussion.

The working vocabulary is:

\begin{longtable}[]{@{}ll@{}}
\toprule\noalign{}
Term & Meaning in this RfD \\
\midrule\noalign{}
\endhead
\bottomrule\noalign{}
\endlastfoot
Stored pixel & One index in the stored image \\
Effective module & A repeated-index 4 × 4 block \\
Dimension & An active hexadecimal palette identifier \\
State & A complete matrix at one temporal index \\
Sequence & An ordered set of up to 64 states \\
Catalogue tuple & A candidate \((x,y,t,d)\) tuple \\
Code & A state-associated value for consumption \\
Key & A dictionary entry or secret proposed for discussion \\
\end{longtable}

English is the default language of the repository and manuscript. A
complete Portuguese edition is maintained alongside it. The
qualifications are collected in \hyperref[scope-and-limitations]{Scope
and Limitations}; the argument below develops the working model
directly.

\section{Working GIF profile and hexadecimal
values}\label{working-gif-profile-and-hexadecimal-values}

\subsection{Geometry and active
entries}\label{geometry-and-active-entries}

The current profile is described by

\[
W_p=H_p=256,\qquad s=4,\qquad W=H=64,\qquad K\le64,\qquad F\le64.
\]

The stored frame has \(W_pH_p=65{,}536\) pixels. Its independent module
count is

\[
WH=\frac{W_pH_p}{s^2}=4{,}096.
\]

The factor \(s^2=16\) repeats each selected index across a square
module. It determines the rendering geometry, while \(WH\) determines
the number of independent module choices.

Each active palette entry is a six-digit hexadecimal value such as
\(\mathtt{\#12ABEF}\). For state \(t\), the mapping is

\[
P_t:\{0,\ldots,K-1\}\longrightarrow
\{\mathtt{\#000000},\ldots,\mathtt{\#FFFFFF}\}.
\]

The public relationship is

\[
\mathrm{dimension}\;\longleftrightarrow\;
\mathrm{palette\ index}\;\longleftrightarrow\;
\mathrm{hexadecimal\ value}.
\]

The current adapter treats the ordered palette as the codebook. Palette
order can change between states while each entry remains a hexadecimal
value. \href{https://www.w3.org/Graphics/GIF/spec-gif89a.txt}{GIF89a}
supplies the indexed table and image-block container.

\subsection{Representation and table
cost}\label{representation-and-table-cost}

A six-digit hexadecimal value has a 24-bit integer domain:

\[
16^6=2^{24}=16{,}777{,}216.
\]

The active alphabet has \(K\) entries, so choosing an entry costs
\(\log_2K\) bits in the unrestricted counting model. At \(K=64\), that
choice carries six bits.

A GIF table with \(K\) entries occupies

\[
B_{\mathrm{table}}=3K\ \text{bytes}.
\]

At 64 entries, the table occupies 192 bytes. If every state supplies its
own table, the table allocation is \(3KF\) bytes. Textual hexadecimal
descriptions and the encoded container have their own storage budgets.
The model keeps these quantities separately observable.

\section{The four-coordinate
catalogue}\label{the-four-coordinate-catalogue}

\subsection{Cartesian catalogue and selected
state}\label{cartesian-catalogue-and-selected-state}

Let

\[
X=\{0,\ldots,W-1\},\quad
Y=\{0,\ldots,H-1\},\quad
T=\{0,\ldots,F-1\},\quad
\Delta=\{0,\ldots,K-1\}.
\]

The proposed four-coordinate catalogue is

\[
\Omega=X\times Y\times T\times\Delta,\qquad
|\Omega|=WHFK.
\]

For the maximum profile,

\[
|\Omega|=64^4=16{,}777{,}216.
\]

A dense sequence chooses a dimension at every position and temporal
state:

\[
D:T\times Y\times X\longrightarrow\Delta.
\]

Equivalently, define a four-index indicator tensor

\[
A[t,y,x,d]=
\begin{cases}
1,&D[t,y,x]=d,\\
0,&\text{otherwise}.
\end{cases}
\]

The selection condition is

\[
\sum_{d=0}^{K-1}A[t,y,x,d]=1.
\]

The catalogue enumerates every candidate tuple; the dense sequence
selects one dimension for each of its \(WHF\) positions. This gives a
precise meaning to the project's four logical coordinates.

\subsection{A hexadecimal catalogue
address}\label{a-hexadecimal-catalogue-address}

With all four coordinate ranges bounded by 64, a row-major catalogue
address can be written as

\[
a=(((64t+y)64+x)64+d).
\]

It spans

\[
0\le a\le64^4-1=\mathtt{0xFFFFFF}.
\]

The inverse is

\[
d=a\bmod64,\qquad
x=\lfloor a/64\rfloor\bmod64,
\]

\[
y=\lfloor a/64^2\rfloor\bmod64,\qquad
t=\lfloor a/64^3\rfloor.
\]

An address uses 24 bits, with six bits for each coordinate. The
implementation exposes the address as a six-digit hexadecimal integer.
This address is useful for discussing dictionaries, sparse events, and
state-local references.

\section{Capacity and nominal rates}\label{capacity-and-nominal-rates}

\subsection{From choices to bits}\label{from-choices-to-bits}

For \(N=WHF\) independently selectable modules, the number of dense
sequences is

\[
|\mathcal D|=K^N.
\]

A binary message of length \(L\) requires \(2^L\) distinguishable
encodings. Therefore

\[
2^L\le K^N
\quad\Longrightarrow\quad
L_{\max}=\lfloor WHF\log_2K\rfloor.
\]

At the maximum profile,

\[
N=64^3=262{,}144,\qquad
\log_2K=6,
\]

\[
L_{\max}=1{,}572{,}864\ \text{bits},
\qquad
B_{\max}=\left\lfloor\frac{L_{\max}}8\right\rfloor
=196{,}608\ \text{bytes}=192\ \text{KiB}.
\]

The per-state gross payload is

\[
B_{\mathrm{state}}=\frac{64^2\cdot6}{8}
=3{,}072\ \text{bytes}=3\ \text{KiB}.
\]

The number of possible complete sequences is exponential in \(N\). The
binary payload is its logarithm and grows proportionally to \(F\).

\subsection{Catalogue rate and payload
rate}\label{catalogue-rate-and-payload-rate}

Let \(f\) be the nominal number of displayed states per second. A full
cycle has duration

\[
\tau=F/f.
\]

For \(F=64\) and \(f=8\), \(\tau=8\) seconds. The catalogue traversal
budget is

\[
R_{\Omega}=\frac{WHFK}{\tau}
=WHKf=2{,}097{,}152
\ \text{candidate tuples/s}.
\]

This is the intended interpretation if the abbreviation \textbf{cps} is
used for the catalogue in this RfD. Two additional rates describe
selected data:

\[
R_{\mathrm{modules}}=WHf=32{,}768
\ \text{modules/s},
\]

\[
R_{\mathrm{bits}}=WHf\log_2K
=196{,}608\ \text{bits/s}
=24\ \text{KiB/s}.
\]

\begin{longtable}[]{@{}lr@{}}
\toprule\noalign{}
Measurement & Maximum profile \\
\midrule\noalign{}
\endhead
\bottomrule\noalign{}
\endlastfoot
Four-coordinate candidates & 16,777,216 per cycle \\
Candidate-tuple budget & 2,097,152 per second at 8 fps \\
Selected modules & 262,144 per cycle \\
Gross payload & 192 KiB per cycle \\
Gross payload rate & 24 KiB/s at 8 fps \\
\end{longtable}

For a future framing design with \(h\) reserved bits per state and code
rate \(r\), a first payload budget is

\[
B_{\mathrm{net}}=
\left\lfloor\frac{rF(WH\log_2K-h)}8\right\rfloor,
\qquad 0<r\le1.
\]

This expression assumes fixed per-state reservations followed by an
aggregate code of rate \(r\). A different placement of reservations
yields a different budget. The eventual layout is a discussion item.

\section{Time in the container and time in the
app}\label{time-in-the-container-and-time-in-the-app}

\subsection{Transport cadence}\label{transport-cadence}

GIF delays are stored in hundredths of a second. A nominal 125 ms state
period can be represented by alternating 120 and 130 ms delays. For an
even frame count,

\[
\sum_{i=0}^{F-1}\delta_i=125F\ \text{ms}.
\]

The prototype quantizes cumulative boundaries:

\[
\delta_i=
10\left(
\left\lfloor\frac{(i+1)p}{10}\right\rfloor
-\left\lfloor\frac{ip}{10}\right\rfloor
\right)\ \text{ms},
\]

where \(p\) is the requested period in milliseconds. At \(p=125\) and
\(F=64\), the stored delays sum to 8,000 ms.

\subsection{App-held sequence and
selection}\label{app-held-sequence-and-selection}

The proposed interaction places the sequence in the app. The user or a
receiving service selects or consumes one state. The app can advance at
a chosen interval \(\Delta\), including one state per minute.

One possible state selector is

\[
i(\theta)=
\left\lfloor\frac{\theta-\theta_0}{\Delta}\right\rfloor\bmod F.
\]

If the user selects at time \(\theta_c\), the frozen index is

\[
i_c=i(\theta_c).
\]

A chosen state-to-code mapping \(g\) supplies the clipboard value:

\[
c=g(D_{i_c}).
\]

The local study viewer uses the state payload's canonical hexadecimal
string for \(g\). It defaults to \(\Delta=60\) seconds. A 64-state cycle
then spans 64 minutes. The 8 fps transport example and the
one-frame-per-minute interaction are two uses of the same ordered
states.

\begin{longtable}[]{@{}
  >{\raggedright\arraybackslash}p{(\linewidth - 4\tabcolsep) * \real{0.3333}}
  >{\raggedright\arraybackslash}p{(\linewidth - 4\tabcolsep) * \real{0.3333}}
  >{\raggedright\arraybackslash}p{(\linewidth - 4\tabcolsep) * \real{0.3333}}@{}}
\toprule\noalign{}
\begin{minipage}[b]{\linewidth}\raggedright
Event
\end{minipage} & \begin{minipage}[b]{\linewidth}\raggedright
App action
\end{minipage} & \begin{minipage}[b]{\linewidth}\raggedright
Receiving side
\end{minipage} \\
\midrule\noalign{}
\endhead
\bottomrule\noalign{}
\endlastfoot
Load & Keep the sequence and its mapping locally & Await a code \\
Advance & Display the next state & Observe or request selection \\
Click & Stop advancement and freeze the displayed state & User selects
the current state \\
Copy & Copy the associated hexadecimal code & User, service, or site
consumes it \\
Resume & Continue when requested & Another state may be selected \\
\end{longtable}

The example renders each selected state as an independent GIF. Clicking
clears the app timer before writing to the clipboard. This makes the
selected image and copied code refer to the same state.

\subsection{Relationship to TOTP}\label{relationship-to-totp}

The comparison with
\href{https://www.rfc-editor.org/rfc/rfc6238.html}{TOTP} concerns the
ownership of evolving state and the role of the receiving verifier. In
TOTP, an app holds a secret and computes a code using a temporal
counter. In this proposal, the app holds or obtains a sequence and
exposes a selected state's code.

The receiving service is the consumer or verifier in both interaction
patterns. The RfD invites discussion of how \(g\), provisioning, state
renewal, and acceptance rules should be chosen for this sequence-based
model. The \href{https://www.w3.org/TR/clipboard-apis/}{Clipboard API}
provides the user-triggered copy operation used by the local example.

\section{Unicode and one-to-one hexadecimal
mappings}\label{unicode-and-one-to-one-hexadecimal-mappings}

\subsection{The scalar domain}\label{the-scalar-domain}

The
\href{https://www.unicode.org/versions/Unicode18.0.0/core-spec/chapter-3/}{Unicode
scalar domain} is

\[
\mathcal U=
\{0,\ldots,\mathtt{0x10FFFF}\}
\setminus
\{\mathtt{0xD800},\ldots,\mathtt{0xDFFF}\}.
\]

Its cardinality is

\[
|\mathcal U|=
\mathtt{0x110000}-\mathtt{0x800}
=1{,}112{,}064.
\]

The proposal's ``approximately one million codes'' is naturally related
to this domain.
\href{https://www.rfc-editor.org/rfc/rfc3629.html}{UTF-8} is a
variable-length encoding of those scalar values using one to four
octets.

A direct hexadecimal representation is

\[
h(u)=\operatorname{hex}_6(u).
\]

It is injective: \(h(u_1)=h(u_2)\) implies \(u_1=u_2\). The
representation fits inside the 24-bit value domain already used by the
catalogue and palette descriptions.

\subsection{Three candidate
interpretations}\label{three-candidate-interpretations}

\textbf{Direct palette dictionary.} A scalar value can be associated
with a six-digit hexadecimal entry. A frame may use at most \(K\)
distinct scalar entries in this interpretation. An adaptive dictionary
could change the active entries between states. Repeated uses of an
entry remain ordinary repeated module choices.

\textbf{Four-dimension scalar representation.} A fixed alphabet of 64
dimension identifiers can represent any scalar with four base-64 digits:

\[
u=64^3d_0+64^2d_1+64d_2+d_3.
\]

The digits are

\[
d_j=
\left\lfloor\frac{u}{64^{3-j}}\right\rfloor\bmod64,
\qquad j=0,1,2,3.
\]

Since

\[
64^3<|\mathcal U|<64^4,
\]

four fixed-length base-64 digits are sufficient, while three cannot
identify every scalar. The prototype provides this reversible mapping.

\textbf{Catalogue-address dictionary.} The scalar's integer can also be
embedded in the 24-bit catalogue address domain. The inverse address
equations produce a coordinate tuple. This supports a one-to-one
dictionary entry or an addressed event. A sequence of addressed events
also needs a representation of order and multiplicity; these become
explicit design choices.

\begin{longtable}[]{@{}
  >{\raggedright\arraybackslash}p{(\linewidth - 6\tabcolsep) * \real{0.2500}}
  >{\raggedright\arraybackslash}p{(\linewidth - 6\tabcolsep) * \real{0.2500}}
  >{\raggedright\arraybackslash}p{(\linewidth - 6\tabcolsep) * \real{0.2500}}
  >{\raggedright\arraybackslash}p{(\linewidth - 6\tabcolsep) * \real{0.2500}}@{}}
\toprule\noalign{}
\begin{minipage}[b]{\linewidth}\raggedright
Scalar example
\end{minipage} & \begin{minipage}[b]{\linewidth}\raggedright
Hexadecimal scalar
\end{minipage} & \begin{minipage}[b]{\linewidth}\raggedright
UTF-8 octets
\end{minipage} & \begin{minipage}[b]{\linewidth}\raggedright
Four dimension identifiers
\end{minipage} \\
\midrule\noalign{}
\endhead
\bottomrule\noalign{}
\endlastfoot
Latin capital A & \texttt{0x000041} & \texttt{41} &
\texttt{00\ 00\ 01\ 01} \\
Latin small a with tilde & \texttt{0x0000E3} & \texttt{C3\ A3} &
\texttt{00\ 00\ 03\ 23} \\
Euro sign & \texttt{0x0020AC} & \texttt{E2\ 82\ AC} &
\texttt{00\ 02\ 02\ 2C} \\
Grinning face & \texttt{0x01F600} & \texttt{F0\ 9F\ 98\ 80} &
\texttt{00\ 1F\ 18\ 00} \\
\end{longtable}

\subsection{Text budget and the eventual
key}\label{text-budget-and-the-eventual-key}

The four-dimension representation uses 24 bits per scalar:

\[
B_{\mathrm{scalar}}=3\ \text{bytes},
\qquad
C_{\mathrm{state}}=\left\lfloor\frac{4096}{4}\right\rfloor=1024
\ \text{scalars}.
\]

A maximum 64-state cycle holds 65,536 such fixed-length scalar
representations before framing. Relative to UTF-8, the scalar mapping
spends more bits on single-octet values and fewer bits on four-octet
values. The useful comparison therefore depends on the text repertoire
and the dictionary design.

The RfD leaves the term ``key'' open to a concrete interpretation: a
public dictionary, a provisioned private dictionary, an index into
external state, or an input to a cryptographic construction. The
one-to-one mapping establishes a representation; its eventual role can
be selected through discussion.

\section{Reproducible case study and app
example}\label{reproducible-case-study-and-app-example}

\subsection{The 16-state study}\label{the-16-state-study}

The published demonstration deliberately contains 16 states. Each state
uses the current 256 × 256 pixel canvas, 64 × 64 effective modules, and
64 hexadecimal palette entries. The ordered palette is rotated between
states. The payload is synthetic and reproducible from seed 20261104.

Its gross data budget is

\[
B_{16}=\frac{64\cdot64\cdot16\cdot6}{8}
=49{,}152\ \text{bytes}=48\ \text{KiB}.
\]

The decoder recovers the original payload exactly. The report records
the original and recovered SHA-256 values, the GIF hash, environment
versions, hexadecimal palettes, state count, and nominal duration.

\begin{longtable}[]{@{}lr@{}}
\toprule\noalign{}
Property & Measured study \\
\midrule\noalign{}
\endhead
\bottomrule\noalign{}
\endlastfoot
Stored image & 256 × 256 pixels \\
Effective state & 64 × 64 modules \\
Frames & 16 \\
Palette entries per state & 64 \\
Stored nominal cycle & 2,000 ms \\
Recovered payload & 49,152 bytes \\
Encoded GIF size & 383,644 bytes \\
Recovery & Byte-exact \\
\end{longtable}

The encoded file size is a measured container size. The gross payload
count comes from the number of independent selected dimensions.

\subsection{Implementation and checks}\label{implementation-and-checks}

The code accepts hexadecimal payload octets and hexadecimal palette
entries. It writes complete indexed frames, including repeated states,
and reads the stored index data through the declared palette order. The
decoder checks that each effective module contains 4 × 4 identical
indices.

The test suite exercises:

\begin{itemize}
\tightlist
\item
  Exact recovery for varying payload lengths, including zero octets.
\item
  Preservation of changing palette order.
\item
  Distinct frame records for identical states.
\item
  The 64-frame, 192 KiB working-profile boundary.
\item
  One state per minute timing.
\item
  Bounds for frames, palettes, and hexadecimal input.
\item
  Truncated-container rejection and inconsistent-module rejection.
\item
  Four-coordinate address inversion and Unicode scalar boundaries.
\end{itemize}

The local viewer defaults to English, offers Portuguese, and pauses
before copying the current state. It includes manual copying when the
direct clipboard operation is unavailable. The complete copied example
code contains 6,144 hexadecimal digits for one 3 KiB state.

The repository provides the case-study generator, results, viewer, both
manuscript sources, and build commands.

\section{Scope and Limitations}\label{scope-and-limitations}

This section consolidates the interpretation and boundaries of the RfD.

\textbf{Discussion status.} The profile, vocabulary, and mappings are
working proposals. Framing, code length, interoperability, provisioning,
and a verification protocol are subjects for discussion before a
specification.

\textbf{Logical dimensions and counting.} The four axes enumerate
candidate tuples. One dimension is selected per spatial-temporal
position. The 16,777,216 catalogue entries and the 192 KiB gross
selected payload measure different objects. ``cps'', when used, denotes
candidate tuples per second.

\textbf{Hexadecimal interface and container.} Hexadecimal is the public
representation of palette values and codes. GIF89a supplies its internal
serialization. A change of numeral representation does not create
additional entropy. The current adapter uses complete opaque frames and
power-of-two palettes with at most 64 entries. The original payload
length is supplied by the caller.

\textbf{Budgets and timing.} Header fields, dictionaries, authentication
data, redundancy, and container storage require separate budgets. The 8
fps example states nominal playback timing; the app's state cadence is
independent. The study does not measure network transfer speed,
generative-model throughput, or physical capture.

\textbf{Authentication.} The app example contains public synthetic
states and demonstrates selection and copying. A repeated finite
sequence needs a chosen provisioning and renewal design for
authentication use. Confidentiality, challenge binding, expiry, replay
handling, and rate limiting belong to that design.
\href{https://www.rfc-editor.org/rfc/rfc6238.html}{TOTP} and
\href{https://www.rfc-editor.org/rfc/rfc5116.html}{AEAD} are references
for the discussion.

\textbf{Unicode and dictionaries.} Scalar values, assigned characters,
grapheme clusters, UTF-8 octets, active palette entries, and keys have
different domains. The scalar mappings preserve scalar identity;
normalization, ordering, multiplicity, and dictionary storage require
explicit choices in a text or key design.

\textbf{Evidence.} The recorded results apply to the current exact
indexed-file path. They support the specified counting and recovery
claims and do not establish the security of a later protocol.

\section{Questions for discussion}\label{questions-for-discussion}

\begin{enumerate}
\def\labelenumi{\arabic{enumi}.}
\tightlist
\item
  Which applications benefit from the four-coordinate catalogue: dense
  data, addressed events, state selection, or dictionary lookup?
\item
  Which hexadecimal object should be called a key, and how should it be
  provisioned?
\item
  Should text use UTF-8 octets, four-dimension scalar values, or
  adaptive state dictionaries?
\item
  How should the selected code be associated with a service, challenge,
  or acceptance window?
\item
  Which applications need fast transport, slow app cadence, or explicit
  user-driven advancement?
\item
  What should renew a sequence after a cycle or after consumption?
\item
  Which fields are necessary before a specification can be proposed?
\end{enumerate}

\section{Authorship and assistance}\label{authorship-and-assistance}

\textbf{João Paulo De Luca Jr and Jorge Luiz Espanhol} are the
co-authors of this RfD and the associated discussion proposal.
ChatGPT/Codex, OpenAI, assisted with mathematical exposition,
translation, implementation, execution, and document preparation. The
authors retain responsibility for the research direction and decisions.

\section{Linked technical references}\label{linked-technical-references}

\begin{enumerate}
\def\labelenumi{\arabic{enumi}.}
\tightlist
\item
  \href{https://www.w3.org/Graphics/GIF/spec-gif89a.txt}{CompuServe,
  Graphics Interchange Format Version 89a}, reproduced by W3C.
\item
  \href{https://pillow.readthedocs.io/en/stable/handbook/image-file-formats.html\#gif}{Pillow,
  GIF image-file support}.
\item
  \href{https://www.unicode.org/versions/Unicode18.0.0/core-spec/chapter-3/}{Unicode
  18.0, Chapter 3: Conformance, especially §3.9}.
\item
  \href{https://www.rfc-editor.org/rfc/rfc3629.html}{RFC 3629: UTF-8, a
  transformation format of ISO 10646}.
\item
  \href{https://www.rfc-editor.org/rfc/rfc6238.html}{RFC 6238: TOTP,
  Time-Based One-Time Password Algorithm}.
\item
  \href{https://www.rfc-editor.org/rfc/rfc5116.html}{RFC 5116: An
  Interface and Algorithms for Authenticated Encryption}.
\item
  \href{https://www.w3.org/TR/clipboard-apis/}{W3C, Clipboard API and
  events}.
\end{enumerate}

\end{document}
